% =====================================================================
% Bagian 3 — Metodologi  (target 2,0 halaman)
% Konfigurasi pelatihan, praproses, dan daftar trade-off dipindahkan ke
% Lampiran A. Definisi metrik ke Lampiran B.
% =====================================================================

\section{Metodologi}

Kontribusi inti beroperasi pada Lapisan~1 (konfirmasi drone) dari arsitektur
berjenjang: Lapisan~0 (pemicu satelit) merupakan integrasi antarmuka pihak
ketiga, dan Lapisan~2 (menara pemantau tetap) tetap berstatus rancangan
konseptual yang tidak divalidasi data uji.

\subsection{Akuisisi data berjenjang dan pemisahan partisi}

Pelatihan berjalan dua tahap, bukan memilih salah satu sumber. Tahap
pra-pelatihan memakai \textbf{53.011 pasangan} RGB--termal FLAME~2
\cite{flame2_2022} berlabel per bingkai, untuk membentuk representasi awal
\emph{encoder} termal yang tidak memiliki \emph{backbone} terlatih berskala
besar. Tahap \emph{fine-tuning} memakai \textbf{anotasi piksel RFFNet}
\cite{rffnet2026} (lisensi MIT) atas subset FLAME~2, dengan partisi resmi
repositori: 552 latih, 240 validasi, 200 uji.

Karena kedua sumber beririsan, pemisahan partisi menjadi gerbang keras. Seluruh
440 indeks bingkai yang menjadi partisi validasi dan uji RFFNet dikecualikan
dari himpunan latih FLAME~2 melalui daftar eksklusi yang diverifikasi otomatis
pada setiap skrip pelatihan. Pencocokan antar-sumber memakai indeks bingkai,
bukan nama berkas, karena skema penamaan kedua sumber berbeda. Praproses dan
augmentasi dirinci pada Lampiran~\ref{app:konfig}.

\subsection{Arsitektur}

Berbeda dengan \emph{concept paper} penyisihan yang menyajikan target rancangan,
seluruh angka pada Tabel~\ref{tab:arsitektur} merupakan nilai aktual yang
dihitung dari \emph{checkpoint} terlatih.

\begin{table}[!htbp]
\caption{Spesifikasi arsitektur dan jumlah parameter aktual.}
\label{tab:arsitektur}
\centering
\small
\begin{tabular}{llrr}
\hline
\textbf{Komponen} & \textbf{Spesifikasi} & \textbf{Parameter} & \textbf{Terlatih} \\
\hline
\emph{Encoder} RGB & DINOv2 ViT-S/14, 9 dari 12 blok dibekukan & 22.056.576 & 5.326.464 \\
\emph{Encoder} termal & CNN 4 blok, dilatih dari nol & 1.273.256 & 1.273.256 \\
Fusi & \emph{Cross-attention} 2 arah, 2 lapis, 6 \emph{head} & 7.100.928 & 7.100.928 \\
\emph{Head} klasifikasi & MLP, 3 kelas & 197.635 & 197.635 \\
\emph{Gate} keandalan & Dua MLP terpisah, keluaran sigmoid & 115.202 & 115.202 \\
\emph{Head} segmentasi & \emph{Decoder} ringan & 856.261 & 856.261 \\
\hline
\textbf{Total} & & \textbf{30.626.467} & \textbf{13.896.355} \\
\hline
\end{tabular}
\end{table}

Fusi bekerja dua arah, RGB melakukan \emph{attend} ke termal dan sebaliknya
pada setiap lapis, sehingga konteks RGB dapat menekan \emph{false positive}
titik panas non-api sementara termal menonjolkan area yang tersamar asap.
Analisis kegagalan pada Subbagian~4.9 memberikan bukti empiris bagi mekanisme
pertama. Beberapa keputusan rancangan mengorbankan sesuatu, resolusi ubin,
pembekuan blok, jumlah lapis fusi, dan luas pencarian $p$, dan
\emph{trade-off}-nya diuraikan pada Lampiran~\ref{app:konfig}.

\subsection{Dari \emph{dropout} pasif ke \emph{gating} adaptif}

Selama pelatihan, salah satu \emph{embedding} modalitas dijatuhkan secara acak
per sampel dengan probabilitas $p$; konfigurasi final memakai $p = 0{,}2$ dari
grid $\{0{,}1;\ 0{,}2;\ 0{,}3\}$. Modul ini aktif hanya pada mode pelatihan.

\emph{Modality dropout} membuat model tahan terhadap kehilangan sensor, tetapi
tidak memberi tahu siapa pun modalitas mana yang sedang diandalkan. Kami
memperluasnya menjadi \emph{gating} adaptif: sebuah \emph{gate} ringan
memprediksi skor keandalan $\in [0,1]$ per modalitas dari token sebelum fusi,
lalu skor tersebut memodulasi bobot \emph{attention} sekaligus mengisi medan
keandalan pada keluaran sistem. Supervisinya diperoleh gratis dari pipeline
augmentasi, target regresi adalah $1 - \tau$, dengan $\tau$ level degradasi
yang memang kami kendalikan sendiri, sehingga tidak ada anotasi tambahan yang
diperlukan. Berbeda dengan protokol kurva degradasi, $\tau_{\text{RGB}}$ dan
$\tau_{\text{termal}}$ di sini disampel \emph{independen} per sampel, karena
\emph{gate} harus menilai tiap modalitas terpisah: asap hanya mengganggu RGB,
\emph{drift} kalibrasi hanya mengganggu termal, dan keduanya jarang terjadi
bersamaan.

\subsection{Dua jalur lokalisasi}

\textbf{Jalur A} memakai \emph{attention rollout} atas bobot \emph{attention}
fusi, tanpa label piksel sama sekali, dengan ambang persentil-90 untuk mengubah
peta atensi menjadi wilayah. \textbf{Jalur B} melatih \emph{head} segmentasi di
atas \emph{encoder} yang dibekukan, memakai mask piksel RFFNet. Jalur A relevan
secara praktis karena lanskap baru umumnya tidak memiliki anotasi piksel.

\subsection{Protokol evaluasi}

Ketahanan diukur sebagai kurva kontinu atas $\tau \in \{0{,}0;\ 0{,}2;\ 0{,}4;\
0{,}6;\ 0{,}8;\ 1{,}0\}$, bukan uji biner fusi-melawan-unimodal. Pada kanal RGB,
degradasi berupa lapisan asap sintetis yang dicampur secara \emph{alpha blend}
dengan opasitas termodulasi derau spasial frekuensi rendah. Pada kanal termal,
degradasi menggabungkan \emph{drift} pola tetap khas sensor mikrobolometer yang
butuh kalibrasi ulang, derau Gaussian, dan pengaburan. Seluruh degradasi
diterapkan pada citra nyata di ruang piksel mentah \emph{sebelum} normalisasi;
kanal termal \textbf{tidak pernah} disintesis dari RGB, konsisten dengan kritik
kami terhadap pendekatan termal sintetis pada Bagian~2.

Konfigurasi unimodal direalisasikan dengan menolkan token modalitas yang
dijatuhkan pada model yang sama, sehingga ketiga kurva berasal dari satu bobot
model, bukan dari tiga model yang dilatih terpisah. Ini mengukur ketahanan
terhadap kehilangan sensor saat inferensi, bukan membandingkan tiga arsitektur.

Karena ukuran partisi kecil, setiap angka utama disertai selang kepercayaan
Wilson 95\%, dan perbandingan antar-konfigurasi memakai \textbf{uji McNemar
berpasangan} dua sisi, bukan perbandingan selang yang tumpang tindih, karena
seluruh konfigurasi dievaluasi pada sampel yang sama. Definisi lengkap seluruh
metrik ada pada Lampiran~\ref{app:metrik}. Karena kegagalan mendeteksi kebakaran
dini jauh lebih mahal daripada alarm palsu, ambang keputusan dipilih pada titik
yang memaksimalkan \emph{recall} dengan presisi tetap di atas batas minimum,
bukan pada titik yang memaksimalkan F1.

\subsection{Sistem terpasang}

Model diintegrasikan ke konsol operator empat halaman dengan kontrak keluaran
tunggal yang disepakati lintas-tim sebelum implementasi dimulai, memuat label
dan keyakinan prediksi, skor keandalan per modalitas, jalur peta lokalisasi
beserta metodenya, sumber pemicu, dan keputusan operator. Lapisan pemicu satelit
menarik data hotspot NASA FIRMS VIIRS SNPP NRT melalui proksi sisi server
(Lampiran~\ref{app:firms}). Alur pengguna ujung-ke-ujung, dari pemicu sampai
keputusan tercatat, beserta penanganan empat kondisi gagal, diuraikan pada
Lampiran~\ref{app:alur}; rancangan antarmuka pada Lampiran~\ref{app:antarmuka};
dan diagram alur eksperimen serta arsitektur sistem, keduanya dipertahankan dari
\emph{concept paper} penyisihan, pada Lampiran~\ref{app:diagram}.
